% cSpell:ignore MSI,SSI,BCD,MUX,PLD,FPLA,PAL,GAL,EPLD,CPLD,FPGA
% ====================== 主文件导言区必须加载的宏包 ======================
% \usepackage{amsmath,amssymb}
% \usepackage{tikz}
% \usepackage{circuitikz}
% \usepackage{caption} % 用于 \captionof 表格小标题
% 本文件为章节子文件，宏包请在主.tex中统一加载，无需重复导入
% ======================================================================

\chapter{组合逻辑电路}
\setlength{\parskip}{0.4em}
\setlength{\itemsep}{0.2em}
\setlength{\tabcolsep}{4pt} % 统一表格列间距

\section{概述}
\subsection{组合逻辑电路的特点}
\begin{itemize}
    \item 任意时刻的输出仅取决于该时刻的输入，与电路原来的状态无关；
    \item 电路结构上信号流向是单向的，没有从输出端到输入端的反馈；
    \item 基本单元是逻辑门电路，不含记忆元件；但门电路有延时，故组合电路也有延迟时间。
\end{itemize}
判断一个电路是否属于组合逻辑电路，就看它有没有记忆元件与反馈通路：有存储单元、输出与历史状态有关的是时序逻辑电路。

\subsection{逻辑功能的描述}
多输入、多输出的组合电路，其输出与输入的逻辑关系可写成一组函数式：
\[
\begin{cases}
y_1=f_1\left(a_1,a_2,\cdots,a_n\right)\\
y_2=f_2\left(a_1,a_2,\cdots,a_n\right)\\
\quad\vdots\\
y_m=f_m\left(a_1,a_2,\cdots,a_n\right)
\end{cases}
\qquad\text{简记为 } Y=F(A)
\]
完整描述一个组合电路可用四种工具：逻辑表达式、真值表、逻辑电路图、时序波形图。四者可相互转换，其中真值表最直观、表达式最便于化简与实现。

\section{组合逻辑电路的分析方法}
分析就是给定逻辑电路图，反推它实现的逻辑功能。标准四步：
\begin{enumerate}
    \item 由逻辑图从输入到输出逐级写出输出端的逻辑表达式；
    \item 用逻辑代数公式或卡诺图化简，得到最简表达式；
    \item 由最简表达式列出真值表；
    \item 由真值表的输入输出规律归纳电路的逻辑功能。
\end{enumerate}

\subsection{分析中常见的功能结论}
\begin{itemize}
    \item 输出为 \(AB+\overline{A}\,\overline{B}\) 或 \(A\overline{B}+\overline{A}B\)：同或、异或功能；
    \item 输出在多输入中"取值不同时为 1、相同时为 0"：判断不一致电路，如三变量的
    \(Y=A\overline{B}+\overline{A}C+B\overline{C}\)；
    \item 输出两路分别为 \(A\oplus B\) 与 \(AB\)：半加器；
    \item 输出在输入 1 的个数过半时为 1：多数表决电路。
\end{itemize}

\subsection{易错点}
\textbf{易错点：}写表达式必须"逐级"进行，先用中间变量表示每一级的输出，再代入化简，不要一步跳到输出端，否则极易漏括号、把与非写成与；化简后必须回到真值表验证功能，只凭表达式猜功能容易出错。

\section{组合逻辑电路的基本设计方法}
设计就是根据实际逻辑问题，求出实现该功能的最简单逻辑电路。"最简"指器件数最少、器件种类最少、连线也最少。

\subsection{设计步骤}
\begin{enumerate}
    \item 分析事件的因果关系，确定输入变量与输出变量；
    \item 定义逻辑状态的含义，即逻辑赋值；
    \item 根据因果关系列出真值表；
    \item 由真值表写出逻辑函数式（最小项之和）；
    \item 用卡诺图化简逻辑函数式；
    \item 按所选器件类型变换函数形式（如统一用与非门），画出逻辑电路图；
    \item 实验验证。
\end{enumerate}

设计过程也可用框图表示：逻辑问题 \(\rightarrow\) 逻辑真值表 \(\rightarrow\) 逻辑函数式 \(\rightarrow\) 选定器件类型 \(\rightarrow\)（SSI：化简后画逻辑图；MSI/PLD：变换后画逻辑图）。

\subsection{逻辑状态的赋值原则}
赋值是人为规定的，但必须遵循两条：一是"1"应赋给使电路正常工作的状态，便于阅读；二是赋值一旦确定，全表必须一致。同一问题赋值方式不同，得到的表达式不同，但逻辑功能等价。

\subsection{函数形式的变换}
实际器件种类有限，常需把最简与或式变换成特定门实现的形式。
\begin{enumerate}
    \item \textbf{与或式 \(\rightarrow\) 与非-与非式}：整体两次取反，再用德摩根定律把内层与项取反。例如三人表决
    \[
    Y=AB+BC+CA=\overline{\overline{AB}\cdot\overline{BC}\cdot\overline{CA}}
    \]
    \item \textbf{与或式 \(\rightarrow\) 或非-或非式}：先求反函数并化为与或式，再整体两次取反。
    \item \textbf{与或式 \(\rightarrow\) 与或非式}：整体取反一次，把结果写成"与项之或再取反"的 \(Y=\overline{AB+CD}\) 形式，可用与或非门一级实现。
\end{enumerate}
用与非门实现时门数通常最少，故工程上最常用。

\section{常用组合逻辑模块}
中规模集成（MSI）组合模块把常用功能做成标准芯片，设计时直接选用并做函数变换，比用 SSI 门搭建更省器件、更可靠。

\subsection{编码器}
\textbf{编码}：为区分一系列不同事物，把每个事物用一组二值代码表示。\textbf{编码器}：把输入的每一个高低电平信号变成一个对应的二进制代码。分普通编码器与优先编码器，又分二进制编码器与二-十进制编码器。

\subsubsection{普通二进制编码器（8 线-3 线）}
输入 \(I_0\sim I_7\) 高电平有效，输出 \(Y_2Y_1Y_0\) 为三位二进制代码。其约束是\textbf{任何时刻只允许一个输入信号有效}，否则输出混乱。由真值表直接写出
\[
Y_2=I_4+I_5+I_6+I_7,\quad
Y_1=I_2+I_3+I_6+I_7,\quad
Y_0=I_1+I_3+I_5+I_7
\]
即每一位输出是"该位为 1 的那些输入"的或。

\subsubsection{优先编码器 74HC148}
允许几个输入同时有效，但设计时已排好优先顺序，只对优先级最高的信号编码。74HC148 是 8 线-3 线优先编码器，输入 \(I_7'\sim I_0'\) 低电平有效、\(I_7'\) 优先权最高，输出 \(Y_2'Y_1'Y_0'\) 为反码，另有三个控制端：
\begin{itemize}
    \item 选通输入 \(S'\)：\(S'=1\) 时电路不工作，所有输出锁定高电平；\(S'=0\) 时正常工作；
    \item 选通输出 \(Y_S'\)：\(S'=0\) 且无编码输入（\(I_7'\sim I_0'\) 全为 1）时 \(Y_S'=0\)，用于多片级联时指示"本片无输入"；
    \item 扩展输出 \(Y_{EX}'\)：\(S'=0\) 且有编码输入时 \(Y_{EX}'=0\)，用于扩展输出位。
\end{itemize}
输出表达式为
\[
Y_2'=\overline{\left(I_4+I_5+I_6+I_7\right)S}
\]
\[
Y_1'=\overline{\left(I_2I_4'I_5'+I_3I_4'I_5'+I_6+I_7\right)S}
\]
\[
Y_0'=\overline{\left(I_1I_2'I_4'I_6'+I_3I_4'I_6'+I_5I_6'+I_7\right)S}
\]
附加输出为
\[
Y_S'=\overline{I_0'I_1'I_2'I_3'I_4'I_5'I_6'I_7'S},\qquad
Y_{EX}'=\overline{\left(I_0+I_1+\cdots+I_7\right)S}
\]
三种状态要记牢：\(S'=1\) 不工作（输出全 1、\(Y_S'=Y_{EX}'=1\)）；\(S'=0\) 但无输入（输出全 1、\(Y_S'=0\)、\(Y_{EX}'=1\)）；\(S'=0\) 且有输入（输出为编码反码、\(Y_S'=1\)、\(Y_{EX}'=0\)）。

\subsubsection{二-十进制优先编码器 74LS147}
把 9 个输入信号 \(I_9'\sim I_1'\) 编成 9 个 BCD 代码，\(I_9'\) 优先权最高、\(I_1'\) 最低，输出 \(Y_3'Y_2'Y_1'Y_0'\) 为 BCD 码的反码。输入全为高电平时输出 1111，相当于对 \(I_0'\) 编码（即数字 0）。

\subsection{译码器}
译码是编码的逆运算：把每个输入的二进制代码译成对应的输出高低电平信号。分二进制译码器、二-十进制译码器、显示译码器。

\subsubsection{二进制译码器与最小项译码器}
把 \(N\) 位二进制代码译成 \(2^N\) 个高低电平信号，称 \(N\) 线-\(2^N\) 线译码器。用二极管与门阵列即可实现，其优点是电路简单，缺点是输入电阻低、输出电阻高、存在输出电平偏移，故只用于集成电路内部。各输出端逻辑式为
\[
Y_i=m_i\qquad(i=0,1,\cdots,2^N-1)
\]
即每个输出对应一个最小项，故称\textbf{最小项译码器}。集成译码器的输出常为低电平有效，此时 \(Y_i'=m_i'\)。

\subsubsection{3 线-8 线译码器 74HC138}
由 CMOS 门构成，地址输入 \(A_2A_1A_0\)，输出 \(Y_7'\sim Y_0'\) 低电平有效，附加控制端 \(S_1\)（高有效）、\(S_2'\)、\(S_3'\)（低有效）。令 \(S=S_1\cdot\overline{S_2'}\cdot\overline{S_3'}\)，则
\[
Y_i'=\overline{S\,m_i}
\]
\begin{itemize}
    \item \(S_1=0\) 或 \(S_2'+S_3'=1\)：译码器禁止工作，输出全为高电平；
    \item \(S_1=1\) 且 \(S_2'=S_3'=0\)：正常译码，被选中端输出低电平，其余为高电平。
\end{itemize}
由于 \(Y_i'=\overline{S m_i}\)，它还可以作\textbf{数据分配器}：把 \(S_2'=S_3'=0\)、数据由 \(S_1\) 输入，则数据被送到 \(A_2A_1A_0\) 指定的那一端，并以反码形式输出。

\subsubsection{二-十进制译码器 74HC42}
把 10 个 BCD 代码译成 10 个高低电平输出，输出逻辑式 \(Y_i'=m_i'\ (i=0\sim9)\)。BCD 码以外的伪码（1010～1111）不产生任何低电平输出。

\subsubsection{显示译码器与数码管}
半导体数码管（LED 七段显示器）把七个发光段排成 8 字形，分共阴极与共阳极两类：共阴极型高电平点亮，共阳极型低电平点亮。液晶显示器（LCD）靠外加电场改变液晶分子排列来透光或遮光，本身不发光、功耗极低。

BCD-七段显示译码器有 4 个输入端、7 个输出端。74LS48 输出高电平有效、配共阴极数码管，74LS47 输出低电平有效、配共阳极数码管。设计方法是：由真值表分别画出 \(Y_a\sim Y_g\) 的卡诺图，圈"0"后求反，得
\[
Y_a=\overline{A_3'A_2'A_1'A_0+A_3A_1+A_2A_0'},\qquad
Y_b=\overline{A_3A_1+A_2A_1A_0'+A_2A_1'A_0}
\]
\[
Y_c=\overline{A_3A_2+A_2'A_1A_0'},\qquad
Y_d=\overline{A_2A_1A_0+A_2A_1'A_0'+A_2'A_1'A_0}
\]
\[
Y_e=\overline{A_2A_1'+A_0},\qquad
Y_f=\overline{A_3'A_2'A_0+A_2'A_1+A_1A_0},\qquad
Y_g=\overline{A_3'A_2'A_1'+A_2A_1A_0}
\]
注意它\textbf{不是最小项译码器}：4 位输入只译出 7 个输出，且伪码也有输出，广义上仍属译码器。

7448 有三个低电平有效的附加控制端：
\begin{enumerate}
    \item 灯测试端 \(LT'\)：\(LT'=0\) 时 \(Y_a\sim Y_g\) 全为 1，数码管显示"8"，用于检查数码管各段是否完好；
    \item 灭零输入 \(RBI'\)：当 \(A_3A_2A_1A_0=0000\) 且 \(RBI'=0\) 时七段全为 0，熄灭该位的 0；
    \item 灭灯输入/灭零输出 \(BI'/RBO'\)：作输入时 \(BI'=0\) 强制熄灭数码管（与输入代码无关）；作输出时仅当 \(A_3A_2A_1A_0=0000\) 且 \(RBI'=0\) 时 \(RBO'=0\)，表示"本位应显示的零已被熄灭"。
\end{enumerate}
利用 \(RBI'\) 与 \(RBO'\) 的配合可实现多位显示的灭零控制：整数部分最高位的 \(RBO'\) 接次高位的 \(RBI'\)，小数部分最低位的 \(RBO'\) 接次低位的 \(RBI'\)，即可自动熄灭无效的前导零与末尾零。

\subsection{数据选择器}
数据选择就是从一组数据中选出某一个送到输出端，也称多路开关。以双 4 选 1 数据选择器 74HC153 为例：数据端 \(D_{10}\sim D_{13}\)，地址端 \(A_1A_0\)，使能端 \(S_1'\) 低有效，输出 \(Y_1\)。当 \(S_1'=0\) 时
\[
Y_1=D_{10}A_1'A_0'+D_{11}A_1'A_0+D_{12}A_1A_0'+D_{13}A_1A_0
\]
即地址码选中哪个数据端，就把该端数据送到输出。

8 选 1 数据选择器 74HC151 的输出为
\[
Y=\sum_{i=0}^{7}m_iD_i
\]
其中 \(m_i\) 是地址 \(A_2A_1A_0\) 的最小项。

\subsubsection{用数据选择器实现逻辑函数}
数据选择器的输出本身就是"地址变量最小项与数据端的与、再或"，因此任何 \(n\) 变量逻辑函数都能用 \(2^n\) 选 1 数据选择器实现：
\begin{enumerate}
    \item 把函数写成最小项之和形式；
    \item 让地址端接 \(n\) 个输入变量；
    \item 对照两式，函数中含某最小项则该数据端接 1，不含则接 0，直接连线即可。
\end{enumerate}
若选用地址数少于变量数的选择器（如三变量函数用 4 选 1），则把多出的一个变量"降维"处理：把它写到数据端，数据端可取 0、1、该变量或其反变量。这是 MSI 实现组合函数最常用的技巧。

\subsection{加法器}
\subsubsection{半加器与全加器}
\textbf{半加器}只考虑两个 1 位二进制数相加，不考虑低位进位：
\[
S=A\overline{B}+\overline{A}B=A\oplus B,\qquad CO=AB
\]
\textbf{全加器}还要考虑低位进位 \(CI\)，实现三个 1 位二进制数相加：
\[
S=A\oplus B\oplus CI,\qquad CO=AB+CI\left(A\oplus B\right)
\]
全加器可由两个半加器加一个或门构成，典型集成器件为双全加器 74LS183。

\subsubsection{多位加法器}
\begin{itemize}
    \item \textbf{串行进位（行波进位）加法器}：低位进位输出直接接高位进位输入，结构简单，但进位逐级传递，运算速度慢，用于速度要求不高的场合；
    \item \textbf{超前进位（快速进位）加法器}：运算开始时即同时得到各位进位。由全加器真值表可知，向高位的进位只可能来自两种情况——\(A_iB_i=1\)，或 \(A_i+B_i=1\) 且 \(CI_i=1\)，故
    \[
    \left(CO\right)_i=A_iB_i+\left(A_i+B_i\right)\left(CI\right)_i
    \]
    令进位生成函数 \(G_i=A_iB_i\)、进位传递函数 \(P_i=A_i+B_i\)，并利用 \((CI)_i=(CO)_{i-1}\) 逐级展开：
    \[
    \left(CO\right)_i=G_i+P_iG_{i-1}+P_iP_{i-1}G_{i-2}+\cdots+P_iP_{i-1}\cdots P_1G_0+P_iP_{i-1}\cdots P_0\left(CI\right)_0
    \]
    第 \(i\) 位和为 \(S_i=A_i\oplus B_i\oplus\left(CI\right)_i\)。
\end{itemize}
74LS283 就是 4 位超前进位加法器：各位和与最后进位基本同时产生，速度快；但电路复杂，且位数越多越复杂。

\subsection{数值比较器}
\subsubsection{1 位数值比较器}
比较 1 位二进制数 \(A\)、\(B\)，三种结果：
\[
Y_{(A>B)}=A\overline{B},\qquad
Y_{(A<B)}=\overline{A}B,\qquad
Y_{(A=B)}=\overline{A\oplus B}=A\odot B
\]

\subsubsection{多位数值比较器}
比较两个多位数必须\textbf{自高位向低位}逐位比较：高位一旦出现大小差异即可判定结果；高位全部相等时才由低位决定。以 4 位比较为例
\[
Y_{(A>B)}=A_3\overline{B_3}+\left(A_3\odot B_3\right)A_2\overline{B_2}
+\left(A_3\odot B_3\right)\left(A_2\odot B_2\right)A_1\overline{B_1}
\]
\[
+\left(A_3\odot B_3\right)\left(A_2\odot B_2\right)\left(A_1\odot B_1\right)A_0\overline{B_0}
+\left(A_3\odot B_3\right)\left(A_2\odot B_2\right)\left(A_1\odot B_1\right)\left(A_0\odot B_0\right)I_{(A>B)}
\]
\(Y_{(A<B)}\) 只需把每项中的原变量与反变量对调；\(Y_{(A=B)}\) 为四位同或之积再与 \(I_{(A=B)}\) 相与。

74LS85（CC14585）为 4 位数值比较器，有 \(A_3\sim A_0\)、\(B_3\sim B_0\) 两组数据输入，\(I_{(A>B)}\)、\(I_{(A<B)}\)、\(I_{(A=B)}\) 三个级联输入，以及三个比较结果输出。使用规则：
\begin{itemize}
    \item 单片比较两个 4 位数时，令 \(I_{(A>B)}=I_{(A<B)}=0\)、\(I_{(A=B)}=1\)；
    \item 多片级联时，低位片的三个输出接到相邻高位片的三个级联输入，高位片即可在"高位相等"时引用低位结果；
    \item 比较时只有高位全部相等的条件下级联输入才被采用，故级联输入对结果无影响的情况在功能表中以"\(\times\)"表示。
\end{itemize}

\section{层次化和模块化的设计方法}
较复杂的组合电路必须采用层次化、模块化的设计方法。
\begin{itemize}
    \item \textbf{层次化}："自顶向下"把总任务分层分块，降低每层复杂度；或"自底向上"先实现并验证有限复杂度的模块，再逐级调用拼装；
    \item \textbf{模块化}：把经过设计与验证的逻辑电路封装成模块反复使用，减少重复劳动、提高可靠性。
\end{itemize}

\subsection{MSI 模块的级联扩展}
用标准模块搭大系统，关键是"用控制端扩展地址/位数"。三种典型扩展：
\begin{enumerate}
    \item \textbf{两片 74HC148 扩展为 16 线-4 线优先编码器}：第一片接高 8 位输入 \(A_{15}'\sim A_8'\)、优先权高，第二片接 \(A_7'\sim A_0'\)；第一片的 \(Y_S'\) 接第二片的 \(S'\)，使第一片有输入时第二片被禁止；第一片的 \(Y_{EX}'\) 经反相后作输出最高位；低 3 位输出取两片输出的"与非"。
    \item \textbf{两片 74HC138 扩展为 4 线-16 线译码器}：用最高位 \(D_3\) 控制两片的使能端——\(D_3=0\) 时第一片工作、第二片禁止，\(D_3=1\) 时相反，从而把 16 个最小项分给两片。
    \item \textbf{两片 CC14585 扩展为 8 位数值比较器}：低位片的三个输出接高位片的三个级联输入，高位片比较高位 4 位、相等时引用低位结果。
\end{enumerate}
扩展的共同思路是：地址/数据的高位用来"选片"，低位用来"片内选址"。

\section{组合逻辑电路中的竞争-冒险}
\subsection{竞争-冒险现象及其成因}
门电路的两个输入信号同时向相反的逻辑电平跳变（一个从 1 变 0、另一个从 0 变 1）的现象称为\textbf{竞争}。

以与门为例：稳态下 \(A=0,B=1\) 与 \(A=1,B=0\) 输出都应为 0。但若 \(A\) 从 1 变 0 的时刻比 \(B\) 从 0 变 1 略早，在 \(A\) 已降到 \(V_{IL(\max)}\) 以下而 \(B\) 尚未升到 \(V_{IH(\min)}\) 以上的短暂 \(\Delta t\) 内，\(A\)、\(B\) 同时处于低电平，输出端就会出现一个极窄的 \(Y=0\) 尖峰脉冲，违背稳态逻辑关系。或门同理，会出现极窄的 \(Y=1\) 尖峰。

应当注意：\textbf{有竞争现象不一定产生尖峰脉冲}。若在 \(B\) 升到 \(V_{IL(\max)}\) 之前 \(A\) 已降到 \(V_{IL(\max)}\) 以下（或门的对应情形），输出端不会出现尖峰。但实际电路中 \(A\)、\(B\) 往往经不同传输途径到达，到达先后与上升下降时间的细微差异难以预知，因此只能说：只要存在竞争，输出就有可能产生违背稳态逻辑的尖峰脉冲。

\textbf{竞争-冒险}：由于竞争而在电路输出端可能产生尖峰脉冲的现象。典型例子是 2 线-4 线译码器：稳态下 \(Y_0\)、\(Y_3\) 都应为 0，但在 \(AB\) 从 10 跳到 01 的过程中，由于两个反相器传输延迟不同，\(Y_0\)、\(Y_3\) 端都会出现尖峰。若负载是对尖峰敏感的电路（如触发器），就会造成误动作，设计时必须设法避免。

\subsection{检查竞争-冒险现象的方法}
在输入变量每次只有一个改变状态的简单情况下，可由逻辑函数式判断。
\begin{enumerate}
    \item \textbf{代数法}：若输出端的逻辑函数在一定条件下能化简成
    \[
    Y=A+\overline{A}\qquad\text{或}\qquad Y=A\overline{A}
    \]
    则存在竞争-冒险。若输出端是或非门、与非门，则应能写成 \(Y=\overline{A+\overline{A}}\) 或 \(Y=\overline{A\overline{A}}\) 的形式。
    \item \textbf{卡诺图法}：若最简与或式对应的卡诺图中存在两个相切（相邻但未被同一圈覆盖）的圈，则存在竞争-冒险，因为相切处对应某个变量取反后两项相接的情形。
\end{enumerate}
代数法的局限是：多数情况下输入变量有两个以上同时改变状态，难以从函数式上简单找出所有竞争-冒险；变量很多时更困难。工程上还可借助计算机模拟检查，或直接用实验观察输出端有无尖峰——加入输入端的信号波形应包含所有可能的状态变化。模拟只能采用标准化典型参数，故\textbf{只有实验检查的结果才是最终结论}。

\subsection{消除竞争-冒险现象的方法}
\begin{enumerate}
    \item \textbf{接入滤波电容}：竞争-冒险产生的尖峰一般很窄（多在几十纳秒以内），在输出端并接一个很小的滤波电容 \(C_f\) 就足以把尖峰幅度削弱到门电路阈值电压以下。TTL 电路 \(C_f\) 通常取几十至几百皮法；CMOS 输出电阻较高，\(C_f\) 可选得更小。优点是简单易行，缺点是增加输出电压波形的上升、下降时间，使波形变坏，故只适用于对波形前后沿无严格要求的场合。
    \item \textbf{引入选通脉冲}：在电路中引入一个选通脉冲 \(p\)，让它的高电平出现在电路达到稳定状态之后，则各门输出端都不会出现尖峰。但此时各门的正常输出也变成脉冲信号，宽度与选通脉冲相同，例如输入 \(AB\) 变成 11 后 \(Y_3\) 并不马上变高，而要等 \(p\) 的正脉冲到来才给出正脉冲。该法不需增加电路元件，但必须得到与输入信号同步的选通脉冲，对脉冲宽度和作用时间有严格要求。
    \item \textbf{修改逻辑设计（增加冗余项）}：例如
    \[
    Y=AB+\overline{A}C=AB+\overline{A}C+BC
    \]
    增加 \(BC\) 项后，在 \(B=C=1\) 时无论 \(A\) 如何改变，输出始终保持 \(Y=1\)，\(A\) 的状态变化不再引起竞争-冒险。\(BC\) 对函数而言是多余的，故称\textbf{冗余项}，此法称增加冗余项法。其适用范围有限：若 \(A\)、\(B\) 同时改变状态（\(AB\) 从 10 变为 01），电路仍存在竞争-冒险，增加的 \(BC\) 只消除了 \(B=C=1\) 时由 \(A\) 变化引起的冒险。
\end{enumerate}
三种方法比较：接滤波电容简单但波形变坏；引选通脉冲不需增加元件但对时序要求严格；修改逻辑设计若运用得当，有时只需增加一根连线（把电路中已存在的某个门输出引到输出门的一个输入端）即可，不必增加门电路。

\section{可编程逻辑器件（PLD）}
\subsection{基本概念}
数字集成电路按功能分通用型与专用型两大类。可编程逻辑器件 PLD 按通用器件生产，其逻辑功能由用户通过对器件编程来设定，兼顾灵活性与硬件性能。

\subsection{PLD 的发展与分类}
发展顺序：FPLA（现场可编程逻辑阵列）、PAL（可编程阵列逻辑）、GAL（通用阵列逻辑）、EPLD（可擦除可编程逻辑器件）、CPLD（复杂可编程逻辑器件）、FPGA（现场可编程门阵列）。

\subsection{FPLA 的结构原理}
FPLA 由可编程与阵列、可编程或阵列和输出缓冲器构成，是组合电路与时序电路的通用结构形式：
\begin{itemize}
    \item 组合型 FPLA：只有与、或阵列，无存储单元，实现组合逻辑函数；
    \item 时序型 FPLA：阵列后增加触发器，可实现时序逻辑电路；
    \item 工作规则：与阵列产生乘积项，或阵列把乘积项或起来得到输出函数；输出缓冲器由 \(\overline{OE}\) 控制，\(\overline{OE}=0\) 时输出使能。
\end{itemize}
LSI 中常用的逻辑图符号包括多输入与阵列、互补输出缓冲器和三态输出缓冲器。

\section{硬件描述语言 Verilog HDL}
\subsection{语言基础}
硬件描述语言用于在不同层次上描述数字电路的结构、功能与行为，工业标准为 Verilog HDL。核心单元是模块 module：
\begin{verbatim}
module <模块名>(<端口列表>);
<内部定义>
<模块条目>
endmodule
\end{verbatim}

\subsection{两种描述方式}
\subsubsection{行为描述方式}
用行为语句描述电路要实现的功能，不关心底层门电路结构。2 选 1 数据选择器：
\begin{verbatim}
module mux_2_to_1 (a,b,out,outbar,sel);
input a,b,sel;
output out,outbar;
assign out = sel ? a : b;
assign outbar = ~out;
endmodule
\end{verbatim}

\subsubsection{结构描述方式}
把硬件电路描述成分级子模块相互连接的结构，调用门原语并用 wire 连线，完全对应硬件原理图：
\begin{verbatim}
module muxgate (a,b,out,outbar,sel);
input a,b,sel;
output out,outbar;
wire out1,out2,selb;
and a1(out1,a,sel);
not i1(selb,sel);
and a2(out2,b,selb);
or  o1(out,out1,out2);
assign outbar = ~out;
endmodule
\end{verbatim}

\subsection{组合电路的描述实例}
\subsubsection{4 位串行进位全加器（结构描述）}
\begin{verbatim}
module Four_bit_fulladd (A,B,CI,S,CO);
parameter size = 4;
input  [size:1] A,B;
input  CI;
output [size:1] S;
output CO;
wire [size-1:1] Ctemp;
onebit_fulladd add1(A[1],B[1],CI,     S[1],Ctemp[1]);
onebit_fulladd add2(A[2],B[2],Ctemp[1],S[2],Ctemp[2]);
onebit_fulladd add3(A[3],B[3],Ctemp[2],S[3],Ctemp[3]);
onebit_fulladd add4(A[4],B[4],Ctemp[3],S[4],CO);
endmodule
\end{verbatim}

\subsubsection{4 位全加器（行为描述）}
\begin{verbatim}
module Four_bit_fulladd (A,B,CI,S,CO);
parameter size = 4;
input  [size:1] A,B;
input  CI;
output [size:1] S;
output CO;
assign {CO,S} = A + B + CI;
endmodule
\end{verbatim}
行为描述用一条加法语句即可得到 5 位结果（最高位为进位），结构描述则要逐个实例化 1 位全加器并连好进位链，二者功能等价。
